\chapter{基于得分的生成模型}
\label{chap:09}
\sourcechapter{9}

% 原书印刷页 217–218，附件 PDF 第 232–233 页。
\section{引言}
\label{sec:09-01}

我们从文献\citeref{1}摘要中的一句话开始。
\begin{quote}
  \emph{从数据产生噪声很容易。从噪声生成数据，就是生成建模。}
\end{quote}

我得说，要给当今的生成建模下一个更简洁的定义，实在不容易。只要看看不同类型的模型，就会马上发现，我们想做的正是从噪声生成数据！不信吗？那就看看各类生成模型怎样工作。
\begin{itemize}
  \item GAN。从已知的 $p(\mathbf z)$ 采样噪声 $\mathbf z$，再用生成器 $G(\mathbf z)$ 得到数据。
  \item VAE。从先验 $p(\mathbf z)$ 采样噪声 $\mathbf z$，再用解码器 $p(\mathbf x\mid\mathbf z)$ 采样数据。
  \item 归一化流。从基分布 $p(\mathbf z)$ 采样噪声 $\mathbf z$，再通过可逆变换 $f$ 得到数据 $\mathbf x=f^{-1}(\mathbf z)$。
\end{itemize}

好奇的读者，希望你已经看出了这里的规律。一般来说，我们在寻找一种变换 $\psi$，将噪声 $\mathbf z$ 映射为数据 $\mathbf x$，即 $\mathbf x=\psi(\mathbf z)$。\figref{fig:09-01} 给出了一个示例。

\dgmfig{fig_09_01}{\linewidth}{从噪声映射到数据的生成建模示意图。通过某种变换 $\psi$，把噪声（左图中的黑色叉号）映射为数据（右图中聚集在橙色月牙附近的黑色叉号）。}{fig:09-01}

在 GAN、VAE 和归一化流中，这个变换由某个神经网络“一步”完成，分别通过生成器、解码器或 $f$ 的逆映射把噪声变成数据。我们也可以设想一种迭代过程，而且已经知道有一类模型正是这么做的，那就是基于扩散的模型！在正向扩散中，每一步都给数据点加入高斯噪声，直到数据点变成纯高斯噪声，即 $\mathbf x_0\sim p_{\mathit{data}}(\mathbf x)$，$\mathbf x_1\sim\mathcal N(\mathbf 0,\mathbf I)$。回顾一下记号，数据点 $\mathbf x\equiv\mathbf x_0$，噪声 $\mathbf z\equiv\mathbf x_1$。接下来，生成部分（也称逆向扩散）从 $\mathbf x_1$（即纯高斯噪声）出发，每一步去掉少量噪声，直到得到数据 $\mathbf x_0$。最初的扩散模型实现\citerefs{2,3}只采用有限步数 $T<+\infty$，所以可以把基于扩散的模型看作具有固定编码器的层次 VAE，参见第~5~章和文献\citeref{4}。

在实践中，扩散模型的成功部分来自其深度。换句话说，我们能够学习极深的层次模型（例如 $T=1000$），因为生成过程被定义为在两个相邻变量之间去掉一小部分噪声。训练时，对 $t\in[0,1]$ 采样，然后优化如下损失，以随机方式训练一个共享网络来预测这微乎其微的噪声。
\begin{equation}
  \mathcal L_t(\mathbf x_t)=\mathbb E_{\mathcal N(\epsilon\mid\mathbf 0,\mathbf I)}
  \bigl[\gamma_t\lVert\epsilon-\epsilon_\theta(\mathbf x_t,t)\rVert^2\bigr],
  \tag{9.1}\label{eq:09-01}
\end{equation}
其中 $\epsilon_t$ 是来自标准高斯分布的样本，$\gamma_t$ 是由噪声调度确定的方差，$\epsilon_\theta(\cdot)$ 是所有 $t$ 共享的神经网络，旨在根据给定的 $\mathbf x_t$ 和 $t$ 预测 $\epsilon$，而 $\mathbf x_t$ 可用给定数据 $\mathbf x_0$ 和 $\epsilon$ 表示。为简单起见，可将 $\gamma_t$ 设为 1，得到所谓的简化目标（simplified objective）\citeref{3}，也可以学习 $\gamma_t$\citeref{5}。

我知道，好奇的读者，你的问题马上就要来了。既然扩散模型效果这么好是因为步数 $T$ 很大，何不索性取无穷多步，即 $T=+\infty$？这个想法合情合理，却带来了一些问题。首先，得处理连续的 $t\in[0,1]$。我们其实已经知道，微分方程可以处理这个问题。于是又出现了第二个问题，怎样通过微分方程定义生成模型？理想情况下，我们想找到一种简单的训练方法，使用类似式 \eqref{eq:09-01} 的损失。事实证明，这可以做到！但我们先来看另一种把噪声变成数据的方法，即得分匹配（score matching）。之后，我们就有了处理连续时间（即无穷深的生成模型）的学习策略。最后还会讨论另一种方法，称为流匹配（flow matching）。

% 原书印刷页 219–225，附件 PDF 第 234–240 页。
\section{得分匹配}
\label{sec:09-02}

\subsection{建模与目标}

请记住，我们的目标是构造一个模型，以迭代的方式把噪声变成数据。下面是第一次尝试。设想有这样一种情形。
\begin{itemize}
  \item 我们能够访问真实分布 $p_{\mathit{real}}(\mathbf x)$，其中 $\mathbf x\in\mathcal X$。
  \item 在数据空间随机取任意一点 $\hat{\mathbf x}\in\mathcal X$，例如从 $\mathcal X$ 上的均匀分布采样。$p_{\mathit{real}}(\hat{\mathbf x})$ 很可能接近 0，即 $p_{\mathit{real}}(\hat{\mathbf x})\approx0$。
  \item \textbf{问题}是，怎样把 $\hat{\mathbf x}$ 变成一个像样的数据点，即现实中可能观测到的对象？
\end{itemize}

实际上，有一种较快、也有理论依据的方法，称为\textbf{随机梯度朗之万动力学}（stochastic gradient Langevin dynamics，SGLD），也称朗之万动力学（Langevin dynamics）\citeref{6}。它可以从 $p_{\mathit{real}}(\mathbf x)$ 采样。从 $\mathbf x_0\equiv\hat{\mathbf x}$ 出发，按以下过程更新。
\begin{equation}
  \mathbf x_{t+\Delta}=\mathbf x_t+\alpha\nabla_{\mathbf x_t}\ln p_{\mathit{real}}(\mathbf x)
    +\eta\cdot\epsilon,
  \tag{9.2}\label{eq:09-02}
\end{equation}
其中 $\Delta$ 是增量，$\alpha>0$，$\eta>0$，且 $\epsilon\sim\mathcal N(\mathbf 0,\mathbf I)$。最后将到达一点 $\mathbf x_0$，满足 $p_{\mathit{real}}(\mathbf x_0)>0$，而且位于某个众数附近。\footnote{译注：原书此处将起点和终点均记为 $\mathbf x_0$。后文第~9.2.3~节的采样框将终点记为 $\mathbf x_1$。}

这很好！不过有一个大问题，我们不知道真实分布 $p_{\mathit{real}}(\mathbf x)$。我们又回到起点，似乎得用 GAN、VAE、归一化流等方法求一个模型 $p_\theta(\mathbf x)$。真要这么做吗？再看式 \eqref{eq:09-02}。我们所需的是实际分布对数的梯度 $\nabla_{\mathbf x}\ln p_{\mathit{real}}(\mathbf x)$。这个量称为\textbf{得分函数}（score function），$\mathbf s(\mathbf x)\overset{\mathrm{df}}{=}\nabla_{\mathbf x}\ln p_{\mathit{real}}(\mathbf x)$。知道 $\mathbf s(\mathbf x)$ 后，就可以运行 SGLD 找到新的点，这些点最终服从真实分布。\figref{fig:09-02} 给出了多峰分布的得分函数示例。\figref{fig:09-02}（a）在网格上用向量表示得分函数。可以看出，得分函数指向各个众数，从而引导朗之万动力学，如\figref{fig:09-02}（b）所示。

\dgmfig{fig_09_02}{.89\linewidth}{（a）得分函数 $\mathbf s(\mathbf x)$ 的示意图。在规则网格上绘制多峰分布的得分向量（白色箭头），颜色越深，概率越低，颜色越亮，概率越高。（b）应用 SGLD 得到的一条轨迹。沿途各点用白色叉号表示，其得分用灰色箭头表示。}{fig:09-02}

因此，\textbf{我们可以拟合得分函数}，从而绕开对数据分布本身的拟合。记住，得分函数是梯度，所以返回的对象与 $\mathbf x$ 形状相同，可以表示为向量。因此，优化如下回归问题，就能找到得分函数的最佳模型\citerefs{7,8}。
\begin{equation}
  \mathcal J(\theta)=\frac12\int
    \lVert\mathbf s_\theta(\mathbf x)-\nabla_{\mathbf x}\ln p_{\mathit{data}}(\mathbf x)\rVert^2
    p_{\mathit{data}}(\mathbf x)\,\mathrm d\mathbf x.
  \tag{9.3}\label{eq:09-03}
\end{equation}
这个积分很常见，就是回归中经常使用的均方误差损失。

式 \eqref{eq:09-03} 的问题是，$p_{\mathit{data}}(\mathbf x)=\frac1N\sum_{n=1}^{N}\delta(\mathbf x-\mathbf x_n)$，其中 $\delta(\cdot)$ 是狄拉克 $\delta$。由于狄拉克 $\delta$ 作为 $\mathbf x$ 的函数不可微，我们无法用自动求导求解式 \eqref{eq:09-03}。解决办法是对数据加上一些方差为 $\sigma^2$ 的小幅高斯噪声，即 $\tilde{\mathbf x}_n=\mathbf x_n+\sigma\cdot\epsilon$，其中 $\epsilon\sim\mathcal N(\mathbf 0,\mathbf I)$。得到的分布是高斯分布 $\mathcal N(\tilde{\mathbf x}_n\mid\mathbf x_n,\sigma^2)$。换句话说，我们可以把 $p_{\mathit{data}}(\mathbf x)$ 变成一个高斯混合模型，以数据点为均值，以较小的 $\sigma^2$ 为方差，即
\begin{equation}
  q_{\mathit{data}}(\tilde{\mathbf x}_n)=\frac1N\sum_{n=1}^{N}
    \mathcal N(\tilde{\mathbf x}_n\mid\mathbf x_n,\sigma^2).
  \tag{9.4}\label{eq:09-04}
\end{equation}

最后，在式 \eqref{eq:09-03} 中用 $q_{\mathit{data}}(\tilde{\mathbf x}_n)$ 替换 $p_{\mathit{data}}(\mathbf x_n)$，就可以构造可微的目标，取代原先的不可微目标\citerefs{7,8,9}。
\begin{equation}
  \begin{aligned}
    \mathcal L(\theta)=\frac{1}{2N}\sum_{n=1}^{N}\int
    &\bigl\lVert\mathbf s_\theta(\tilde{\mathbf x})
        -\nabla_{\tilde{\mathbf x}}\ln\mathcal N(\tilde{\mathbf x}\mid\mathbf x_n,\sigma^2)
      \bigr\rVert^2\\[-2pt]
    &{}\cdot\mathcal N(\tilde{\mathbf x}_n\mid\mathbf x_n,\sigma^2)\,\mathrm d\tilde{\mathbf x}.
  \end{aligned}
  \tag{9.5}\label{eq:09-05}
\end{equation}
也许你会说，还有一个问题，因为这里要算一个梯度。不过，高斯分布的得分函数有闭式表达式。
\begin{align}
  \nabla_{\tilde{\mathbf x}}\ln\mathcal N(\tilde{\mathbf x}_n\mid\mathbf x_n,\sigma^2)
    &=\nabla_{\tilde{\mathbf x}}
      \left(-\ln(2\pi\sigma^D)-\frac{1}{2\sigma^2}(\tilde{\mathbf x}_n-\mathbf x_n)^2\right)
      \tag{9.6}\label{eq:09-06}\\
    &=-\nabla_{\tilde{\mathbf x}}\frac{1}{2\sigma^2}(\tilde{\mathbf x}_n-\mathbf x_n)^2
      \tag{9.7}\label{eq:09-07}\\
    &=-\frac{1}{\sigma^2}(\tilde{\mathbf x}_n-\mathbf x_n)
      \tag{9.8}\label{eq:09-08}\\
    &=-\frac{1}{\sigma^2}(\mathbf x_n+\sigma\cdot\epsilon-\mathbf x_n)
      \tag{9.9}\label{eq:09-09}\\
    &=-\frac{1}{\sigma}\epsilon,
      \tag{9.10}\label{eq:09-10}
\end{align}
其中第四个等式用到了 $\tilde{\mathbf x}$ 的定义 $\tilde{\mathbf x}_n=\mathbf x_n+\sigma\cdot\epsilon$。因此，得分为
\begin{equation}
  \nabla_{\tilde{\mathbf x}}\ln\mathcal N(\tilde{\mathbf x}_n\mid\mathbf x_n,\sigma^2)
  =-\frac{1}{\sigma}\epsilon,
  \tag{5}\label{eq:09-original-5}
\end{equation}
将式 \eqref{eq:09-original-5} 代入式 \eqref{eq:09-05}，得到可微目标的最终形式。
\begin{align}
  \mathcal L(\theta)
    &=\frac{1}{2N}\sum_{n=1}^{N}\int
      \left\lVert\mathbf s_\theta(\tilde{\mathbf x})+\frac{1}{\sigma}\epsilon\right\rVert^2
      \mathcal N(\epsilon\mid\mathbf 0,\sigma^2\mathbf I)\,\mathrm d\epsilon
      \tag{9.11}\label{eq:09-11}\\
    &=\frac{1}{2N}\sum_{n=1}^{N}\mathbb E_{\mathcal N(\epsilon\mid\mathbf 0,\mathbf I)}
      \left[\left\lVert\mathbf s_\theta(\tilde{\mathbf x})+\frac{1}{\sigma}\epsilon\right\rVert^2\right]
      \tag{9.12}\label{eq:09-12}\\
    &=\frac{1}{N}\sum_{n=1}^{N}\mathbb E_{\mathcal N(\epsilon\mid\mathbf 0,\mathbf I)}
      \left[\frac{1}{2\sigma}\left\lVert\epsilon-\tilde{\mathbf s}_\theta(\tilde{\mathbf x})\right\rVert^2\right],
      \tag{9.13}\label{eq:09-13}
\end{align}
最后一行改用了得分模型 $\tilde{\mathbf s}_\theta(\mathbf x)=-\sigma\mathbf s_\theta(\tilde{\mathbf x})$，并利用了 $\sigma$ 为正的性质（别忘了，标准差总是正的！）。\footnote{译注：原书式 \eqref{eq:09-11} 的高斯协方差与式 \eqref{eq:09-12} 及前面定义的标准高斯 $\epsilon$ 不一致。按本段所给的重参数化，式 \eqref{eq:09-13} 从式 \eqref{eq:09-12} 提出的因子应为 $1/(2\sigma^2)$。公式及后续算法和代码保留印刷写法。} 实现时，当然可以直接用神经网络参数化 $\tilde{\mathbf s}_\theta(\tilde{\mathbf x})$，省去先参数化 $\mathbf s_\theta(\mathbf x)$ 再缩放的步骤。只要清楚这里推导的每一步即可。通过优化式 \eqref{eq:09-13} 的目标来学习得分模型 $\tilde{\mathbf s}_\theta(\tilde{\mathbf x})$，称为（去噪）\textbf{得分匹配}（denoising score matching）\citerefs{7,8,10}。

得分匹配让我们以一种全新的方式学习生成模型。通过求解回归问题，可以把\textbf{匹配分布}的任务转成\textbf{匹配得分}的任务。这给生成模型的训练带来了新的视角。但要记住，最终目的是应用朗之万动力学，从噪声生成数据。下面汇总得分匹配在实践中所需的信息。

\subsection{训练}

可以重复以下步骤来学习得分模型。
\begin{booknote}[colframe=blue!65!black,colback=blue!4]{得分匹配的训练}
  \begin{enumerate}
    \item 选取一个数据点 $\mathbf x$。
    \item 从 $\mathcal N(\epsilon\mid\mathbf 0,\mathbf I)$ 采样 $\epsilon$。
    \item 计算加噪后的数据 $\tilde{\mathbf x}=\mathbf x+\sigma\cdot\epsilon$。
    \item 计算得分 $\tilde{\mathbf s}_\theta(\tilde{\mathbf x})$。
    \item 计算得分匹配目标关于 $\theta$ 的梯度，即
      $\Delta_\theta=\nabla_\theta\frac{1}{2\sigma}
      \lVert\epsilon-\tilde{\mathbf s}_\theta(\tilde{\mathbf x})\rVert^2$。
      我们用自动微分完成这一步。
    \item 更新得分模型 $\theta:=\theta-\Delta_\theta$。返回第 1 步，直至满足停止条件。
  \end{enumerate}
\end{booknote}
实际应用中会采用小批量，但整体过程相同。

\subsection{采样（生成）}

学到得分模型后，就可以借助朗之万动力学进行采样。生成过程如下。
\begin{booknote}[colframe=blue!65!black,colback=blue!4]{用朗之万动力学从得分匹配模型采样}
  \begin{enumerate}
    \item 随机采样一个点 $\mathbf x_0$，例如从均匀分布采样。
    \item 运行朗之万动力学 $T$ 步，其中 $\Delta=1/T$。
      \begin{equation}
        \mathbf x_{t+\Delta}=\mathbf x_t+\alpha\tilde{\mathbf s}_\theta(\mathbf x_t)+\eta\cdot\epsilon.
        \tag{9.14}\label{eq:09-14}
      \end{equation}
    \item 返回最终的点 $\mathbf x_1$。
  \end{enumerate}
\end{booknote}

最终的点应服从数据分布。请注意，为保持朗之万动力学的形式，我们这里用的是 $\mathbf s_\theta(\mathbf x_t)$，而非 $\tilde{\mathbf s}_\theta(\mathbf x_t)$。\footnote{译注：这里的原文说明与式 \eqref{eq:09-14} 不一致，式中印为 $\tilde{\mathbf s}_\theta(\mathbf x_t)$。} 这不成问题，因为已知两者的关系为 $\tilde{\mathbf s}_\theta(\mathbf x_t)=-\sigma\mathbf s_\theta(\mathbf x_t)$，所以 $\mathbf s_\theta(\mathbf x_t)=-\frac1\sigma\tilde{\mathbf s}_\theta(\mathbf x_t)$。

\subsection{得分匹配与基于扩散的模型}

好奇的读者，我敢说你已经看出了一些端倪。看看式 \eqref{eq:09-03} 中基于扩散的模型的目标，以及式 \eqref{eq:09-13} 中得分模型的目标，可以清楚看到二者几乎完全相同，尤其是取 $\gamma_t=1/\sigma$ 时！\footnote{译注：这里的式号 \eqref{eq:09-03} 沿用原书引用，基于扩散的模型的目标见式 \eqref{eq:09-01}。} 实际上，我们可以把得分匹配变成一种（近似的）基于扩散的方法。需要做三件事。
\begin{enumerate}
  \item 引入方差 $\sigma_t^2$ 的调度，其中
    \[
      t\in\mathcal T=\{t\colon t\text{ 从 }0\text{ 开始，按增量 }\Delta=1/T\text{ 递增到 }1\}.
    \]
  \item 修改得分模型，让它能够感知这一调度，即 $\tilde{\mathbf s}_\theta(\mathbf x,\sigma_t^2)$。
  \item 修改目标，使它也能感知这一调度。
    \begin{equation}
      \mathcal L(\theta)=\frac{1}{2N}\sum_{n=1}^{N}\sum_{t\in\mathcal T}
      \lambda_t\,\mathbb E_{\mathcal N(\epsilon\mid\mathbf 0,\mathbf I)}
      \left[\frac{1}{\sigma}\left\lVert\epsilon-\tilde{\mathbf s}_\theta(\tilde{\mathbf x},\sigma_t^2)\right\rVert^2\right],
      \tag{8}\label{eq:09-original-8}
    \end{equation}
    其中 $\lambda_t$ 为加权系数，例如 $\lambda_t=\sigma_t^2$。
\end{enumerate}

对此类模型采样需要做些调整。例如，可以采用退火朗之万动力学（annealed Langevin dynamics），见文献\citeref{10}的算法 1。最终得到的采样过程与基于扩散的模型非常相似。差别在于采样方案针对每个 $\sigma_t^2$ 执行多步朗之万动力学，而基于扩散的模型每次“去噪”只执行一步。

\subsection{用代码实现得分匹配}

我们已经概述了如何用得分匹配实现训练和采样（生成）。代码相当直接，唯一的技巧是把数据变换到 $[-1,1]$。下面的代码相比上面的分析只改了一点，使用 $\frac{1}{2\sigma}\lVert\epsilon+\tilde{\mathbf s}_\theta(\tilde{\mathbf x})\rVert^2$，使朗之万动力学仍采用标准形式 $\mathbf x_{t+\Delta}=\mathbf x_t+\alpha\tilde{\mathbf s}_\theta(\mathbf x_t)+\eta\cdot\epsilon$。

\bookcode{code/ch09_01_score_matching.py}{得分匹配类。}{lst:09-01}

得分模型采用多层感知机（MLP），超参数设为 $\alpha=0.1$、$\sigma=0.1$、$\eta=0.05$、$T=100$。运行代码后，可得到类似\figref{fig:09-03} 的结果。

\dgmfig{fig_09_03}{\linewidth}{得分匹配的结果。（a）真实图像的样本。（b）生成图像的样本。（c）验证集上计算的得分匹配损失示例。}{fig:09-03}[（c）中横轴 epochs 为训练轮次，纵轴 score matching loss 为得分匹配损失。]

\subsection{得分匹配能用来做什么？}

\textbf{改进得分匹配。} 可以看到，极其简单的代码就能得到一些相当像样的生成结果。不过，原论文\citeref{10}指出，得分匹配还有以下几个问题。
\begin{itemize}
  \item 首先也是最重要的一点是，低密度区域（很少有数据或没有数据的区域）会给得分估计带来困难。再看\figref{fig:09-02}（a）中颜色较暗区域的真实得分。在这些区域里，几乎看不到任何数据点。因此，得分函数的估计可能很不准确。
  \item 第二个问题由第一个引起。得分模型训练得差，朗之万动力学就可能生成质量较差的样本。在\figref{fig:09-03}（b）的示例中可以多少看出这一点。
  \item 第三，如果真实数据分布的各个众数相隔很远，朗之万动力学就要走许多步才能到达这些位置。这是所谓的\emph{混合缓慢}（slow mixing）问题。
\end{itemize}

为了克服这些问题，可以采用如下方法。
\begin{itemize}
  \item 可以采用一系列方差，并使用前面讨论过的退火朗之万动力学。详细做法请参阅文献\citeref{10}以及介绍了其他实用技巧的后续论文\citeref{11}。
  \item 采用\emph{切片得分匹配}（sliced score matching）\citeref{12}，可以训练出更好的得分模型，尤其是在高维情形下。它的思想类似于切片 Wasserstein 距离（sliced Wasserstein distance）\citeref{13}，把得分函数和得分模型投影到某一随机方向。这样，新的切片得分匹配目标更容易计算，只需做向量乘法。这个新目标与 Hutchinson 迹估计量（Hutchinson's trace estimator）密切相关。
\end{itemize}

\textbf{其他扩展。} 得分匹配是一个通用框架，可以进一步扩展。下面是一些例子。
\begin{itemize}
  \item 得分匹配适用于基于能量的模型。这一思想已用于训练玻尔兹曼机\citerefs{14,15}和其他基于能量的模型\citeref{16}。
  \item 匹配随机变量分布得分的思想也可以用来学习离散随机变量的分布\citeref{17}。为此，需要采用一种适合概率质量函数的得分定义。
  \item 得分匹配的应用可以超出图像范围。文献\citeref{18}提出了一种点云建模方法，采用了适于处理点云的特定得分模型。
\end{itemize}

% 原书印刷页 226–237，附件 PDF 第 241–252 页。
\section[用随机／常微分方程定义生成模型]{以随机微分方程／常微分方程定义生成模型}
\label{sec:09-03}

\subsection{回顾基于扩散的模型}

再来想一想基于扩散的模型。此前，我们将其介绍为层次 VAE，其变分后验有一种特定的构造方式，先对前面的潜变量做线性变换，然后加上高斯噪声。也可以把这类模型看成时间离散化的动力系统。为说明这一点，设我们把数据点 $\mathbf x_0$（$\mathbf x_0\sim p_0(\mathbf x)\equiv p_{\mathit{data}}(\mathbf x)$）变成从已知分布 $\pi$ 采样的噪声 $\mathbf x_1$（$\mathbf x_1\sim p_1(\mathbf x)\equiv\pi(\mathbf x)$）。用 $t\in[0,1]$ 表示“时间”，$\mathbf x_0$ 和 $\mathbf x_1$ 之间有 $T$ 步，步长为 $\Delta=1/T$。于是，\emph{正向扩散}可以表示为
\begin{equation}
  \mathbf x_{t+\Delta}=\sqrt{1-\beta_t}\,\mathbf x_t+\sqrt{\beta_t}\,\epsilon_t,
  \tag{9.15}\label{eq:09-15}
\end{equation}
其中 $\epsilon_t\sim\mathcal N(\mathbf 0,\mathbf I)$。好奇的读者，你可以看出，这是一个时间离散化的动力系统。请记住这一点，稍后我们还会回到这个问题。

基于扩散的模型有一个有趣的性质。由于具有线性和高斯性质，可以从数据 $\mathbf x_0$ 直接计算 $\mathbf x_t$。
\begin{equation}
  \mathbf x_t=\sqrt{\alpha_t}\,\mathbf x_0+\sqrt{1-\alpha_t}\,\epsilon_t,
  \tag{9.16}\label{eq:09-16}
\end{equation}
其中 $\alpha_t=\prod_{\tau=1}^{t}(1-\beta_\tau)$。进一步，可以按如下方式反算数据点。
\begin{equation}
  \mathbf x_0=\frac{1}{\sqrt{\alpha_t}}\mathbf x_t-
    \frac{\sqrt{1-\alpha_t}}{\sqrt{\alpha_t}}\epsilon_t.
  \tag{9.17}\label{eq:09-17}
\end{equation}
正向扩散时，可以采样 $\epsilon_t$。从 $\mathbf x_t$ 逆推 $\mathbf x_0$ 时，却无法访问噪声 $\epsilon_t$。通常是这样，正向扩散一直进行到 $\mathbf x_t$，此前的所有 $\mathbf x$ 和 $\epsilon$ 都不再保留。还有办法！毕竟，我们可以引入神经网络 $\epsilon_\theta(\mathbf x_t,t)$ 来预测噪声。

怎样学习这个神经网络？在奠基论文\citeref{3}中，损失是以下各项损失的和。
\begin{equation}
  \mathcal L_t(\theta)=\lVert\epsilon_\theta(\mathbf x_t,t)-\epsilon_t\rVert^2,
  \tag{9.18}\label{eq:09-18}
\end{equation}
论文表明，它等价于 ELBO。不过请等一等，好奇的读者，你看到我看到的东西了吗？这个损失还等价于另一种损失！没错，就是上一节介绍的得分匹配。得分模型和噪声模型之间有如下对应关系（对标准差为 $\sigma$ 的高斯去噪分布）。
\begin{equation}
  \mathbf s_\theta(\mathbf x,t)=-\frac{\epsilon_\theta(\mathbf x_t,t)}{\sigma}.
  \tag{9.19}\label{eq:09-19}
\end{equation}

把前面的观察归纳如下。
\begin{itemize}
  \item 正向扩散定义了一个离散时间动力系统。
  \item 基于扩散的模型的损失函数与得分匹配损失（几乎）相同。
  \item 基于扩散的模型实际上与得分模型相对应。
  \item 基于扩散的模型与采用 $\sigma$ 调度训练的得分模型非常相似，因为两者都以迭代方式运行，且每一步考虑的噪声都更多。
\end{itemize}
这些相似之处表明，也许存在一个能推广得分模型和扩散模型的通用框架。

\subsection{寻求通用框架}

在解答得分模型与基于扩散的模型能否共用一个框架之前，先回到最初的那句话。“也可以把基于扩散的模型看成时间离散化的动力系统”，因为正向扩散具有如下形式（我们再写一遍，重复一次也无妨）。
\begin{equation}
  \mathbf x_{t+\Delta}=\sqrt{1-\beta_t}\,\mathbf x_t+\sqrt{\beta_t}\,\epsilon_t.
  \tag{9.20}\label{eq:09-20}
\end{equation}

好奇的读者，我们可以坦诚相告。有人没有学过动力系统，有人学过却已经忘了。在着手写本节之前，我对常微分方程（ordinary differential equation，ODE）和随机微分方程（stochastic differential equation，SDE）的记忆已经非常、非常生疏了。因此，我们先简短地探访一下微分方程的世界！对于更好奇的读者（没错，说的就是你！），我强烈推荐文献\citeref{19}这本书，基础和进阶内容兼顾。

先透露一点后面的内容。事实证明，SDE 和 ODE 在生成建模中大有用处。所以，请跟我继续往下读，我们将发现一类新的生成模型，以及它所带来的崭新而美妙的世界！

\subsubsection{常微分方程和数值方法}

我们从 ODE 开始。一般来说，可以将它定义如下。
\begin{equation}
  \frac{\mathrm d\mathbf x_t}{\mathrm dt}=f(\mathbf x_t,t),
  \tag{9.21}\label{eq:09-21}
\end{equation}
给定某个初始条件 $\mathbf x_0$。有时将 $f(\mathbf x_t,t)$ 称为\emph{向量场}（vector field）。

可以用某种数值方法求解这个 ODE，以特定方式对时间进行离散化。例如，欧拉法（Euler's method）如下进行，从 $t=0$ 出发，以 $\Delta$ 为步长推进到 $t=1$。
\begin{align}
  \mathbf x_{t+\Delta}-\mathbf x_t&=f(\mathbf x_t,t)\cdot\Delta
    \tag{9.22}\label{eq:09-22}\\
  \mathbf x_{t+\Delta}&=\mathbf x_t+f(\mathbf x_t,t)\cdot\Delta.
    \tag{9.23}\label{eq:09-23}
\end{align}
有时需要从 $t=1$ 回到 $t=0$，这时可以用反向欧拉法（backward Euler's method）。
\begin{equation}
  \mathbf x_t=\mathbf x_{t+\Delta}-f(\mathbf x_{t+\Delta},t+\Delta)\cdot\Delta.
  \tag{9.24}\label{eq:09-24}
\end{equation}
聪明的读者，你一定看出了其中的联系。如果把 $\mathbf x_0$ 当作数据，把 $\mathbf x_1$ 当作噪声，那么，\emph{只要}知道 $f(\mathbf x_t,t)$，就可以运行反向欧拉法，得到一个生成模型！请先记下这个想法。没错，你想到的正是这样！

\subsubsection{随机微分方程与概率流常微分方程}

现在来看 SDE。一般而言，可以把 SDE 想象成轨迹为随机的 ODE，在每个 $t$，状态按某种概率分布 $p_t(\mathbf x)$ 分布。SDE 可以定义为
\begin{equation}
  \mathrm d\mathbf x_t=f(\mathbf x_t,t)\,\mathrm dt+g(t)\,\mathrm d\mathbf v_t,
  \tag{9.25}\label{eq:09-25}
\end{equation}
其中 $\mathbf v$ 是标准维纳过程（standard Wiener process）。在这一语境中，$f(\cdot,t)$ 称为\emph{漂移}（drift），$g(\cdot)$ 是称为\emph{扩散}（diffusion）的标量函数。漂移分量是确定的（可参照上面的 ODE 写法），扩散分量由于标准维纳过程而具有随机性。关于维纳过程的更多内容，请参阅文献\citeref{19}。在这里，只需知道它具有高斯性质即可。例如，$\mathbf v_{t+\Delta}-\mathbf v_t$ 服从正态分布，$\mathbf v_{t+\Delta}-\mathbf v_t\sim\mathcal N(0,\Delta)$。顺带一提，这里假定漂移和扩散的形式都是\textbf{已知的}。后面会给出一个例子，好奇的读者，现在只要记住它们是给定的。

这个 SDE 有一个重要性质，存在一个对应的常微分方程，其解服从相同的分布！如果从数据点 $\mathbf x_0$ 出发，求解以下概率流常微分方程（probability flow ODE，PF-ODE）\citeref{1}，就能得到噪声 $\mathbf x_1\sim p_1(\mathbf x)$。
\begin{equation}
  \frac{\mathrm d\mathbf x_t}{\mathrm dt}
  =\left(f(\mathbf x_t,t)-\frac12g^2(t)\nabla_{\mathbf x_t}\ln p_t(\mathbf x_t)\right).
  \tag{9.26}\label{eq:09-26}
\end{equation}
再仔细看看这个结果。
\begin{itemize}
  \item 首先，维纳过程不再出现。因此，现在处理的是 ODE，随机项已经消去。
  \item 其次，漂移和扩散分量仍在，但扩散分量平方后乘以 $-\frac12$。
  \item 第三，出现了得分函数！要看推导，请参阅文献\citeref{1}。这里直接采用这个结果。它看起来也合情合理。SDE 的解服从 $p_t(\mathbf x)$，因此得分函数出现在这里并不奇怪。毕竟，得分函数表明按 $p_t(\mathbf x)$ 形成的轨迹应当是什么样子。
\end{itemize}

\subsubsection{将概率流常微分方程用作基于得分的生成模型}

等等！这里得到了什么？暂且假设得分函数已知，就可以从 $\mathbf x_1\sim\pi(\mathbf x)$ 出发，应用反向欧拉法，把这个 PF-ODE 用作生成模型。
\begin{equation}
  \mathbf x_t=\mathbf x_{t+\Delta}-
  \left(f(\mathbf x_{t+\Delta},t+\Delta)-
    \frac12g^2(t+\Delta)\nabla_{\mathbf x_{t+\Delta}}\ln p_t(\mathbf x_{t+\Delta})\right)\cdot\Delta.
  \tag{9.27}\label{eq:09-27}
\end{equation}
太好了！现在的问题是得分函数，但我们已经知道怎样处理。可以用去噪得分匹配来学习它，也就是用以数据点为中心的高斯混合模型作为加噪后的经验分布，再进行得分匹配。与前面介绍的去噪得分匹配相比，这里还需要考虑时间 $t$。
\begin{equation}
  \mathcal L(\theta)=\int_0^1\mathcal L_t(\theta)\,\mathrm dt.
  \tag{9.28}\label{eq:09-28}
\end{equation}
怎样定义 $\mathcal L_t(\theta)$？按得分匹配的思想，应考虑 $\lambda_t\lVert\mathbf s_\theta(\mathbf x_t,t)-\nabla_{\mathbf x_t}\ln p_t(\mathbf x_t)\rVert^2$。但我们算不出 $p_t(\mathbf x_t)$，所以要另想办法。可以定义分布 $p_{0t}(\mathbf x_t\mid\mathbf x_0)$，用于对原始数据点 $\mathbf x_0$ 的加噪版本采样。综合起来，得到
\begin{equation}
  \begin{aligned}
    \mathcal L_t(\theta)={}&\tfrac12\mathbb E_{\mathbf x_0\sim p_{\mathit{data}}(\mathbf x)}
      \mathbb E_{\mathbf x_t\sim p_{0t}(\mathbf x_t\mid\mathbf x_0)}\\
    &\left[\lambda_t\lVert\mathbf s_\theta(\mathbf x_t,t)
      -\nabla_{\mathbf x_t}\ln p_{0t}(\mathbf x_t\mid\mathbf x_0)\rVert^2\right].
  \end{aligned}
  \tag{9.29}\label{eq:09-29}
\end{equation}
重要的是，如果 $p_{0t}(\mathbf x_t\mid\mathbf x_0)$ 取高斯分布，就可以解析计算得分函数。马上会看到一个例子。此外，计算 $\mathcal L_t(\theta)$ 时可以只用一个样本，即取蒙特卡洛估计。

找到 $\mathbf s_\theta(\mathbf x_t,t)$ 后，就可以运行反向欧拉法来采样数据。
\begin{equation}
  \mathbf x_t=\mathbf x_{t+\Delta}-
  \left(f(\mathbf x_{t+\Delta},t+\Delta)-
    \frac12g^2(t+\Delta)\mathbf s_\theta(\mathbf x_{t+\Delta},t+\Delta)\right)\cdot\Delta.
  \tag{9.30}\label{eq:09-30}
\end{equation}
请记住，这里假定漂移和扩散都是已知的。此外，我们沿用（反向）欧拉法，但你尽管选择其他 ODE 求解器，放开手脚试一试！这里我们希望尽量讲得清楚简单。

\figref{fig:09-04} 展示了用反向欧拉法采样的例子。对于某个 PF-ODE 和给定的得分函数，可以从一个多峰分布得到样本。可以想象，ODE 求解器会“走向”众数。这个简单的例子表明，定义 PF-ODE 是一种有力的生成手段！只要得分函数近似得好，就能直接从原始分布采样。

\dgmfig{fig_09_04}{.48\linewidth}{对多峰分布运行反向欧拉法的示例。这里取 $f(\mathbf x,t)=0$、$g(t)=9^t$、$T=100$。得分函数（函数本身，无需模型！）通过自动求导计算。绿色方块表示 $\mathbf x_1\sim\pi$，蓝色圆点是样本 $\mathbf x_0$。}{fig:09-04}

\subsection[基于得分的生成模型示例]{基于得分的生成模型示例：方差爆炸型概率流常微分方程}

\subsubsection{模型的构造}

要定义自己的基于得分的生成模型（score-based generative model，SBGM），需要以下要素。
\begin{itemize}
  \item 漂移 $f(\mathbf x,t)$。
  \item 扩散 $g(t)$。
  \item $p_{0t}(\mathbf x_t\mid\mathbf x_0)$ 的形式。
\end{itemize}
在文献\citerefs{1,20}中可以找到三种 SBGM 示例，分别是方差爆炸型（variance exploding，VE）SDE、方差保持型（variance preserving，VP）SDE 和 sub-VP SDE。这里重点讨论 VE SDE，其漂移和扩散取为
\begin{itemize}
  \item $f(\mathbf x,t)=0$。
  \item $g(t)=\sigma^t$，其中 $\sigma>0$ 是超参数。注意，$\sigma$ 的指数是时间 $t\in[0,1]$。
\end{itemize}
根据上面的讨论，把所选的 $f(\mathbf x,t)$ 和 $g(t)$ 代入 PF-ODE 的一般形式，得到
\begin{equation}
  \frac{\mathrm d\mathbf x_t}{\mathrm dt}
    =-\frac12\sigma^{2t}\nabla_{\mathbf x_t}\ln p_t(\mathbf x_t).
  \tag{9.31}\label{eq:09-31}
\end{equation}

接下来，要学习得分模型，需要定义产生 $\mathbf x_0$ 的加噪版本所用的条件分布。好在 SDE 理论（例如文献\citeref{19}的第 5 章）告诉我们如何计算 $p_{0t}(\mathbf x_t\mid\mathbf x_0)$！具体公式见文献\citeref{1}的附录。这里给出最终的解。
\begin{equation}
  p_{0t}(\mathbf x_t\mid\mathbf x_0)=
  \mathcal N\left(\mathbf x_t\mid\mathbf x_0,
    \frac{1}{2\ln\sigma}(\sigma^{2t}-1)\mathbf I\right).
  \tag{9.32}\label{eq:09-32}
\end{equation}
所以，方差随时间变化的函数为
\begin{equation}
  \sigma_t^2=\frac{1}{2\ln\sigma}(\sigma^{2t}-1).
  \tag{9.33}\label{eq:09-33}
\end{equation}
当 $\sigma$ 足够大时，最终分布 $p_{01}(\mathbf x)$ 接近以下高斯分布。
\begin{align}
  p_{01}(\mathbf x)
    &=\int p_0(\mathbf x_0)*\mathcal N\left(\mathbf x\mid\mathbf x_0,
      \frac{1}{2\ln\sigma}(\sigma^2-1)\mathbf I\right)
      \tag{9.34}\label{eq:09-34}\\
    &\approx\mathcal N\left(\mathbf x\mid\mathbf 0,
      \frac{1}{2\ln\sigma}(\sigma^2-1)\mathbf I\right).
      \tag{9.35}\label{eq:09-35}
\end{align}
我们会用 $p_{01}(\mathbf x)$ 采样噪声 $\mathbf x_1$，再将其逆推为数据 $\mathbf x_0$。\footnote{译注：式 \eqref{eq:09-34} 保留原书写法。该边缘化积分应以 $p_0(\mathbf x_0)\mathcal N(\mathbf x\mid\mathbf x_0,\frac{\sigma^2-1}{2\ln\sigma}\mathbf I)$ 为被积函数，对 $\mathbf x_0$ 积分。}

\dgmfig{fig_09_05}{.49\linewidth}{不同 $\sigma$ 取值下，$p_{0t}(\mathbf x_t\mid\mathbf x_0)$ 的标准差随 $t$ 的变化。}{fig:09-05}[原书图注称“标准差”，但纵轴标的是方差 $\sigma_t^2$。]

\begin{booknote}{旁注}
为了更好地理解 $\sigma$ 的各种取值如何影响 $\sigma_t$，可以看看\figref{fig:09-05} 中的曲线。显然不能取 $\sigma=1$（你问为什么？因为 $\ln(1)=0$，会遇到除以 0 的问题），不过，取 $\sigma=1.01$ 时，得到的 $p_1(\mathbf x)$ 接近标准高斯分布。
\end{booknote}

\subsubsection{怎样选择 \texorpdfstring{$\lambda_t$}{λₜ}}

进入训练过程之前，最后说一下如何选择 $\mathcal L_t(\theta)$ 定义中的 $\lambda_t$。到目前为止我一直略过这一点，是有原因的。Ho 等人\citeref{3}直接取 $\lambda_t\equiv1$。这就完了！真是这样吗？聪明的读者，你也许想得到，事情没那么简单。Song 和 Kingma\citeref{21}表明，对于 VE PF-ODE，设 $\lambda_t=\sigma_t^2$ 实际上有好处。甚至可以用各项 $\mathcal L_t(\theta)$ 的和作为对数似然函数的代理量。我们会利用这一点在训练中实施提前停止。

\subsubsection{训练}

下面以所选的 VE SBGM 为例介绍训练过程。正如前面讨论得分匹配时所说，过程相对简单直接，包含以下步骤。
\Needspace{11\baselineskip}
\begin{booknote}[colframe=blue!65!black,colback=blue!4]{VE SBGM 的训练过程}
  \begin{enumerate}
    \item 选取一个数据点 $\mathbf x_0$。
    \item 采样 $\mathbf x_1\sim\pi(\mathbf x)=\mathcal N(\mathbf x\mid\mathbf 0,\mathbf I)$。
    \item 采样 $t\sim\operatorname{Uniform}(0,1)$。
    \item 计算 $\mathbf x_t=\mathbf x_0+\sqrt{\frac{1}{2\ln\sigma}(\sigma^{2t}-1)}\cdot\mathbf x_1$。这是从 $p_{0t}(\mathbf x_t\mid\mathbf x_0)$ 抽取的样本。
    \item 在 $\mathbf x_t$ 和 $t$ 处计算得分模型 $\mathbf s_\theta(\mathbf x_t,t)$ 的值。
    \item 计算单个样本的得分匹配损失 $\mathcal L_t(\theta)=\sigma_t^2\lVert\mathbf x_1-\sigma_t\mathbf s_\theta(\mathbf x_t,t)\rVert^2$。
    \item 利用 $\nabla_\theta\mathcal L_t(\theta)$，用基于梯度的方法更新 $\theta$。
  \end{enumerate}
\end{booknote}
针对可用的训练数据重复这七步，直到满足某个停止条件。显然，实际训练采用小批量，逐步处理成批的数据点。

在这一训练过程中，我们有意使用 $-\sigma_t\mathbf s_\theta(\mathbf x_t,t)$，因为 $-\sigma_t\mathbf s_\theta(\mathbf x_t,t)=\epsilon_\theta(\mathbf x_t,t)$，所以准则 $\sigma_t^2\lVert\mathbf x_1-\epsilon_\theta(\mathbf x_t,t)\rVert^2$ 与基于扩散的模型相对应\citerefs{4,5}。\footnote{译注：若按本段的得分和噪声关系代入最后的平方范数，训练框第 6 步中的减号应为加号。原框和正文各自的写法均保留。} 现在你明白为什么我们坚持要把基于扩散的模型看作动力系统了吧！

\subsubsection{采样}

得分模型训练完成后，终于可以生成了！为此，需要运行反向欧拉法（别忘了，也可以用其他 ODE 求解器）。对于 VE PF-ODE，它采取如下形式。
\begin{align}
  \mathbf x_t
    &=\mathbf x_{t+\Delta}+\left(\frac12\sigma^{2(t+\Delta)}
      \left\{-\frac{1}{\sigma^{t+\Delta}}\mathbf s_\theta(\mathbf x_{t+\Delta},t+\Delta)\right\}\right)
      \cdot\Delta
      \tag{9.36}\label{eq:09-36}\\
    &=\mathbf x_{t+\Delta}-\left(\frac12\sigma^{t+\Delta}
      \mathbf s_\theta(\mathbf x_{t+\Delta},t+\Delta)\right)\cdot\Delta,
      \tag{9.37}\label{eq:09-37}
\end{align}
起点为 $\mathbf x_1\sim p_{01}(\mathbf x)=\mathcal N\bigl(\mathbf x\mid\mathbf 0,\frac{1}{2\ln\sigma}(\sigma^2-1)\mathbf I\bigr)$。注意，第一行用的是加号，因为 VE PF-ODE 的扩散项是 $-\frac12\sigma^{2t}$，因此反向欧拉法中的减号变成了加号。读者也许觉得这十分明显，可我总会搞混加号和减号，所以宁愿在这里写得明确一点。

\subsection{终于到了代码！}

实现 VE SBGM 的条件已经齐备！与得分匹配类似，我们把数据变换到 $[-1,1]$。你会注意到，这段代码与得分匹配的代码有不少相似之处。这在意料之中，因为两者在概念上非常相似。

\Needspace{6\baselineskip}
\bookcode{code/ch09_02_ve_pf_ode.py}{VE PF-ODE 模型类。}{lst:09-02}

得分模型采用 MLP，超参数设为 $\sigma=1.01$ 和 $T=20$。运行代码后，可以得到与\figref{fig:09-06} 类似的结果。

我很少评论结果，尤其是因为我所用的数据集过于简单（这是有意为之，让你能很快运行代码！）。不过，我想指出两点。
\begin{itemize}
  \item VE SBGM 用起来并不容易，因为噪声相当大。查看生成图像，就能对可能的问题有大致的认识。这里让 ODE 求解器运行 20 步，但 SBGM 通常需要数百步，甚至数千步。
  \item 看看代理量在验证集上的取值。众所周知，SBGM 的收敛速度较慢。
\end{itemize}

\dgmfig{fig_09_06}{\linewidth}{VE PF-ODE 模型的结果。（a）真实图像的样本。（b）生成图像的样本。（c）验证集上计算的得分匹配损失示例。}{fig:09-06}[（c）中横轴 epochs 为训练轮次，纵轴 proxy 为代理量。]

\subsection{基于得分的生成模型还有许多精彩之处！}

\textbf{其他 SBGM。} 除了方差爆炸型，还有其他类型的 SBGM\citeref{1}。
\begin{itemize}
  \item 方差保持型（VP）。漂移为 $f(\mathbf x,t)=-\frac12\beta_t\mathbf x$，扩散为 $g(t)=\sqrt{\beta_t}$，损失权重为 $\lambda_t=1-\exp\{-\int_0^t\beta_s\,\mathrm ds\}$，其中 $\beta_t$ 是时间 $t$ 的某个函数。
  \item Sub-VP。漂移为 $f(\mathbf x,t)=-\frac12\beta_t\mathbf x$，扩散和损失权重分别为
  \[
    \begin{aligned}
      g(t)&=\sqrt{\beta_t\left(1-\exp\left\{-2\int_0^t\beta_s\,\mathrm ds\right\}\right)},\\
      \lambda_t&=\left(1-\exp\left\{-\int_0^t\beta_s\,\mathrm ds\right\}\right)^2,
    \end{aligned}
  \]
  其中 $\beta_t$ 是时间 $t$ 的某个函数。
\end{itemize}
这些模型有多种版本，尤其是 $\lambda_t$ 和其他依赖于 $t$ 的函数可以有不同的定义方式，例如 VE 的 $\sigma_t$ 和 VP 的 $\beta_t$。综述见文献\citeref{5}。

\textbf{更好的求解器。} 前面简要提到过，SBGM 的一个缺点是采样步数（即 ODE 求解器的步数）很多。文献\citeref{22}介绍了专门的 ODE 求解器，在 $T=10$ 步内就能达到很好的性能。文献\citeref{23}又把步数降到了 $T=5$！通常，选用更合适的求解器可以得到更好的结果，例如采用 Heun 法（Heun's method）\citeref{24}。

\textbf{其他改进。} 基于得分的模型领域还有很多想法！这里只举几例。
\begin{itemize}
  \item 在潜在空间中使用 SBGM\citeref{25}。
  \item 实际上，可以用类似神经常微分方程（neural ODE）的方法计算 SBGM 的对数似然函数\citeref{26}。Song 等人\citeref{20}表明，适当修改得分匹配损失，可以得到对数似然函数的上界。
  \item 用去量化（dequantization）、重要性加权（importance weighting）等多种技巧改进对数似然的估计\citeref{27}。
  \item 文献\citeref{28}提出了一种新模型，称为\emph{一致性模型}（consistency model）。其思想是学习一个能一步把噪声匹配到数据的模型。
  \item 文献\citeref{29}提出了把 SBGM 扩展到黎曼流形的方法。
\end{itemize}

\textbf{已经有了大量工作！} 有关 SBGM 的论文还在源源不断地发表。可以在这个网页查看关于 SBGM 的最新概况。
\begin{center}
  \url{https://scorebasedgenerativemodeling.github.io/}
\end{center}

% 原书印刷页 238–254，附件 PDF 第 253–269 页。
\section{流匹配}
\label{sec:09-04}

\subsection[从常微分方程看连续归一化流]{从另一视角看基于常微分方程的生成模型：连续归一化流（CNF）}

\subsubsection{再谈常微分方程}

前面讨论了如何通过随机微分方程（SDE），或者等价地通过相应的概率流常微分方程（PF-ODE），定义生成模型。我们说明了用反向欧拉法之类的数值方法求解 SDE 或 ODE，可以得到一个将噪声迭代地转换为数据的生成过程。不过，有一个问题，我们是否必须先写出 SDE 及其等价的 PF-ODE，或者能否用任意 ODE 来定义生成模型？好奇的读者，我知道你的意思！之所以用 SDE，是因为我们知道可以用得分匹配来训练。那么，若取任意 ODE，该怎样学习这类生成模型呢？

我们说得具体一些。回顾 ODE 的定义。
\begin{equation}
  \frac{\mathrm d\mathbf x_t}{\mathrm dt}=v(\mathbf x_t,t),
  \tag{9.38}\label{eq:09-38}
\end{equation}
其中，\emph{向量场} $v(\mathbf x_t,t)$ 定义了动力学。用权重为 $\theta$ 的神经网络 $v_\theta(\mathbf x_t,t)$ 参数化向量场，得到所谓的\textbf{神经常微分方程}\citeref{26}。将这个神经 ODE 的初始条件（例如噪声）记为 $\mathbf x_0$，求解方程，也就是沿时间 $t$ 积分，便得到输出（例如数据）。\footnote{译注：按式 \eqref{eq:09-38} 的初值问题，积分给出 $\mathbf x_1-\mathbf x_0$。式 \eqref{eq:09-39} 保留印刷形式。}
\begin{equation}
  \mathbf x_1=\int_0^1 v_\theta(\mathbf x_t,t)\,\mathrm dt.
  \tag{9.39}\label{eq:09-39}
\end{equation}
到目前为止一切顺利，但与基于得分的生成模型相比，（几乎）没多出什么新东西。后者匹配的是得分，分布匹配则是让模型拟合经验分布。我们能否构造一种基于似然的训练方法？简短的回答是能。容易吗？简短的回答又是，不容易。下面就详细看看这两个答案。

\subsubsection{从连续性方程（质量守恒）到瞬时变量变换}

再说一遍，只要有向量场的模型，从 ODE 采样，也就是从 $t=0$ 积分到 $t=1$，就不难。可以用欧拉法，也可以用其他数值求解器。但是，如果希望把 $v_\theta(\mathbf x,t)$ 诱导出的分布拟合到数据分布，怎样获得这样的模型就是个问题。毕竟，从已知分布 $\mathbf x_0\sim\pi(\mathbf x)$（如标准高斯分布）出发，求解 ODE 会得到另一个分布！我们可以用\textbf{连续性方程}（continuity equation）写出这个诱导分布的解析形式。接下来要用一点数学，甚至还要用点物理，准备好，我们开始吧。

暂且把概率想象成质量，是一种手指可以触到的东西。我总想到黏土，你也可以想成水。现在设想一根管道，沿长度方向的横截面积处处相同，我们的物质（例如水）在其中流动。每一时刻都有一定的质量通量（flux）$f_t$，也就是说，（概率）质量按照向量场（或速度）移动，$f_t(\mathbf x_t)=p_t(\mathbf x_t)v(\mathbf x_t,t)$。既然说的是概率质量，或沿等截面管道流动的水，质量是守恒的。没有新的质量（水）出现或消失，既没有漏出去，也没有倒进来。用数学语言说，质量的变化率 $\frac{\partial p_t(\mathbf x_t)}{\partial t}$ 加上所有方向上通量的变化率（也称通量的散度），其和保持为零，即质量守恒。
\begin{equation}
  \frac{\partial p_t(\mathbf x_t)}{\partial t}
    +\operatorname{div}\bigl(p_t(\mathbf x_t)v(\mathbf x_t,t)\bigr)=0,
  \tag{9.40}\label{eq:09-40}
\end{equation}
其中 $\operatorname{div}(\cdot)$ 是散度（divergence），定义为
\[
  \operatorname{div}(V(x_1,\ldots,x_D))
    =\sum_{d=1}^{D}\frac{\partial V_d(x_1,\ldots,x_D)}{\partial x_d},
\]
即 $V$ 的各分量分别对相应变量求一阶偏导数，再求和。

\begin{booknote}{旁注}
向量场的散度等于雅可比矩阵的迹！对于二维空间中的向量场 $V(x_1,x_2)$，有
\[
  \operatorname{div}(V(x_1,x_2))=
  \frac{\partial V_1(x_1,x_2)}{\partial x_1}+
  \frac{\partial V_2(x_1,x_2)}{\partial x_2}
  =\operatorname{Tr}\left(\frac{\partial V_i(x_1,x_2)}{\partial x_i}\right),
\]
其中
\[
  \frac{\partial V_i(x_1,x_2)}{\partial x_i}=
  \begin{bmatrix}
    \frac{\partial V_1}{\partial x_1}&\frac{\partial V_1}{\partial x_2}\\
    \frac{\partial V_2}{\partial x_1}&\frac{\partial V_2}{\partial x_2}
  \end{bmatrix}
\]
是雅可比矩阵，$\operatorname{Tr}(\mathbf A)=\sum_i A_{ii}$ 是矩阵 $\mathbf A$ 的迹。
\end{booknote}

利用向量微积分的恒等式和散度的性质，可以用概率分布的对数来写连续性方程，这也称为\textbf{瞬时变量变换}（instantaneous change of variables）\citeref{26}。
\begin{equation}
  \frac{\mathrm d\ln p(\mathbf x_t)}{\mathrm dt}+
  \operatorname{Tr}\left(\frac{\partial v(\mathbf x_t,t)}{\partial\mathbf x_t}\right)=0.
  \tag{9.41}\label{eq:09-41}
\end{equation}

\begin{booknote}{旁注}
我知道，好奇的读者，你不喜欢别人只告诉你结论，却不给出证明。这是旁注，如果没时间可以直接跳过。不过，我觉得这个结果挺有意思，也有助于后面的讨论。那我们需要什么呢？首先，有一个一般性质，$\frac{\partial\ln p_t}{\partial t}=\frac{1}{p_t}\frac{\partial p_t}{\partial t}$，或者 $p_t\frac{\partial\ln p_t}{\partial t}=\frac{\partial p_t}{\partial t}$。其次，$\operatorname{div}(p_tv)=\langle\nabla_{\mathbf x}p_t,v\rangle+p_t\operatorname{div}(v)$，所以 $\frac{\operatorname{div}(p_tv)}{p_t}=\langle\nabla_{\mathbf x}\ln p_t,v\rangle+\operatorname{div}(v)$。第三，已经看到向量场的散度等于其雅可比矩阵的迹。现在把这三点代入连续性方程。
\begin{gather}
  \frac{\partial p_t(\mathbf x_t)}{\partial t}+
  \operatorname{div}\bigl(p_t(\mathbf x_t)v(\mathbf x_t,t)\bigr)=0
  \tag{9.42}\label{eq:09-42}\\
  \frac{1}{p_t(\mathbf x_t)}\frac{\partial p_t(\mathbf x_t)}{\partial t}+
  \frac{1}{p_t(\mathbf x_t)}\operatorname{div}\bigl(p_t(\mathbf x_t)v(\mathbf x_t,t)\bigr)=0
  \tag{9.43}\label{eq:09-43}\\
  \frac{\partial\ln p_t(\mathbf x_t)}{\partial t}+
  \langle\nabla_{\mathbf x}\ln p_t,v\rangle+\operatorname{div}\bigl(v(\mathbf x_t,t)\bigr)=0
  \tag{9.44}\label{eq:09-44}\\
  \langle\nabla_{\mathbf x}\ln p_t,v\rangle=
  -\frac{\partial\ln p_t(\mathbf x_t)}{\partial t}-\operatorname{div}\bigl(v(\mathbf x_t,t)\bigr).
  \tag{9.45}\label{eq:09-45}
\end{gather}
接下来，计算 $\ln p_t(\mathbf x_t)$ 的全导数，再把上面求得的 $\langle\nabla_{\mathbf x}\ln p_t,v\rangle$ 代入，就得到\citeref{30}
\begin{align}
  \frac{\mathrm d\ln p(\mathbf x_t)}{\mathrm dt}
  &=\left\langle\frac{\partial\ln p(\mathbf x_t)}{\partial t},
    \frac{\partial t}{\partial t}\right\rangle+
    \left\langle\nabla_{\mathbf x}\ln p(\mathbf x_t),
    \frac{\partial\mathbf x_t}{\partial t}\right\rangle
    \tag{9.46}\label{eq:09-46}\\
  &=\frac{\partial\ln p(\mathbf x_t)}{\partial t}+
    \left\langle\nabla_{\mathbf x}\ln p(\mathbf x_t),
    \frac{\partial\mathbf x_t}{\partial t}\right\rangle
    \tag{9.47}\label{eq:09-47}\\
  &=\frac{\partial\ln p(\mathbf x_t)}{\partial t}+
    \bigl\langle\nabla_{\mathbf x}\ln p(\mathbf x_t),v(\mathbf x_t,t)\bigr\rangle
    \tag{9.48}\label{eq:09-48}\\
  &=\frac{\partial\ln p(\mathbf x_t)}{\partial t}-
    \frac{\partial\ln p_t(\mathbf x_t)}{\partial t}
    -\operatorname{div}\bigl(v(\mathbf x_t,t)\bigr)
    \tag{9.49}\label{eq:09-49}\\
  &=-\operatorname{div}\bigl(v(\mathbf x_t,t)\bigr)
    \tag{9.50}\label{eq:09-50}\\
  &=-\operatorname{Tr}\left(\frac{\partial v(\mathbf x_t,t)}{\partial\mathbf x_t}\right).
    \tag{9.51}\label{eq:09-51}
\end{align}
\end{booknote}

随后沿时间积分，便可计算对数密度的总变化。
\begin{gather}
  \int_0^1\left(\frac{\mathrm d\ln p(\mathbf x_t)}{\mathrm dt}+
    \operatorname{Tr}\left(\frac{\partial v(\mathbf x_t,t)}{\partial\mathbf x_t}\right)\right)
    \,\mathrm dt=0
    \tag{9.52}\label{eq:09-52}\\
  \ln p(\mathbf x_1)-\ln\pi(\mathbf x_0)+
    \int_0^1\operatorname{Tr}\left(\frac{\partial v(\mathbf x_t,t)}{\partial\mathbf x_t}\right)
    \,\mathrm dt=0
    \tag{9.53}\label{eq:09-53}\\
  \ln p(\mathbf x_1)=\ln\pi(\mathbf x_0)-
    \int_0^1\operatorname{Tr}\left(\frac{\partial v(\mathbf x_t,t)}{\partial\mathbf x_t}\right)
    \,\mathrm dt.
    \tag{9.54}\label{eq:09-54}
\end{gather}

为什么费心把所有量都表示成对数概率？因为最后一行正是归一化流所用变量变换的连续版本！这里对雅可比矩阵的迹沿时间积分，离散版本则将雅可比矩阵的对数行列式求和。因此，训练神经 ODE 与训练归一化流相似，只是这里的时间连续。所以，这一语境中的神经 ODE 称为\textbf{连续归一化流}（continuous normalizing flows，CNF）。

\subsubsection{计算连续归一化流的对数似然}

这里和离散时间归一化流有一点不同，我们并不要求 $v$ \emph{可逆}。因此，给定数据点 $\mathbf x_1$，通常不能简单地将变换求逆来得到 $\mathbf x_0$。不过，在相当宽松的条件下，即 $v$ 及其一阶导数均为 Lipschitz 连续（例如使用 SELU、SiLU 等 Lipschitz 连续激活函数的神经网络），就可以唯一地求解下面的问题\citeref{31}。
\begin{equation}
  \begin{bmatrix}
    \mathbf x_0\\
    \ln p_1(\mathbf x_1)-\ln\pi(\mathbf x_0)
  \end{bmatrix}=
  \int_1^0
  \begin{bmatrix}
    v_\theta(\mathbf x_t,t)\\
    -\operatorname{Tr}\left(\dfrac{\partial v_\theta(\mathbf x_t,t)}{\partial\mathbf x_t}\right)
  \end{bmatrix}\,\mathrm dt,
  \tag{9.55}\label{eq:09-55}
\end{equation}
其初始条件如下。
\begin{equation}
  \begin{bmatrix}
    \mathbf x_1\\
    \ln p_1(\mathbf x_{\mathit{data}})-\ln p_1(\mathbf x_1)
  \end{bmatrix}=
  \begin{bmatrix}
    \mathbf x_{\mathit{data}}\\
    0
  \end{bmatrix},
  \tag{9.56}\label{eq:09-56}
\end{equation}
其中 $\mathbf x_1$ 是数据点 $\mathbf x_{\mathit{data}}$，对数概率之差为零。注意，我们是按反向次序求解的，即从数据 $\mathbf x_1$ 回到噪声 $\mathbf x_0$。\footnote{译注：按式 \eqref{eq:09-38} 和式 \eqref{eq:09-54}，式 \eqref{eq:09-55} 右侧的反向积分分别给出 $\mathbf x_0-\mathbf x_1$ 和 $\ln\pi(\mathbf x_0)-\ln p_1(\mathbf x_1)$。左侧保留原书印刷形式。}

\begin{booknote}[colframe=blue!65!black,colback=blue!4]{计算对数似然函数}
归纳起来，需要完成以下步骤。
\begin{enumerate}
  \item 取一个数据点 $\mathbf x_1=\mathbf x_{\mathit{data}}$。
  \item 用求解器求解式 \eqref{eq:09-55}，得到 $\mathbf x_0$ 并记录随时间变化的迹。
  \item 把 $\ln\pi(\mathbf x_0)$ 与负的迹之和
    $-\int_0^1\operatorname{Tr}\bigl(\frac{\partial v_\theta(\mathbf x_t,t)}{\partial\mathbf x_t}\bigr)\,\mathrm dt$
    相加，算出对数似然。
\end{enumerate}
\end{booknote}

现在，我们可以通过求解器反向传播，事情就成了！真是这样吗？整个过程的复杂度如何？乍一看，开销似乎很大。

\subsubsection{Hutchinson 迹估计量}

好奇的读者，如果你记得，归一化流的问题在于计算大小为 $D\times D$ 的雅可比矩阵的对数行列式。一般情形下，成本为 $O(D^3)$。计算迹需要 $O(D^2)$，因为要把对角线元素相加，而每个元素都需要一次单独的前向传播，因此复杂度为二次。比离散情形好一些，不过还有一个问题，我们必须通过数值求解器反向传播！恐怕天下没有免费的午餐。Chen 等人\citeref{26}提出使用\emph{伴随灵敏度法}（adjoint sensitivity method），可以将它看作连续时间的反向传播。随后，通过最大化对数似然 $\ln p(\mathbf x_1)$ 来训练神经 ODE。然而，每来一个新数据点，都需要先用数值方法求解 ODE，再通过求解过程反向传播。这项操作的成本很高！

因此，需要尽可能寻找降低成本的办法。其中一个技巧涉及如何计算迹。利用\textbf{Hutchinson 迹估计量}\citeref{31}，可以把二次复杂度降到 $O(D)$。对于任意方阵 $\mathbf A$，这个估计量相对容易计算。
\begin{equation}
  \operatorname{Tr}(\mathbf A)=\mathbb E_\epsilon[\epsilon^\top\mathbf A\epsilon],
  \tag{9.57}\label{eq:09-57}
\end{equation}
其中 $\epsilon$ 服从均值为零、方差为一的分布，例如 $\epsilon\sim\mathcal N(\mathbf 0,\mathbf I)$。对某个特定的 $\epsilon$，乘积 $\mathbf A\epsilon$ 可通过一次前向传播计算，也可以对其反向传播。因此，可以取 $M$ 个蒙特卡洛样本来估计迹。
\begin{equation}
  \operatorname{Tr}(\mathbf A)\approx\frac1M\sum_{m=1}^{M}\epsilon_m^\top\mathbf A\epsilon_m.
  \tag{9.58}\label{eq:09-58}
\end{equation}
在实践中取 $M=1$，即每来一个新数据点，只取一个 $\epsilon$ 样本。这个估计量显然带有噪声，不过它是无偏的。所以，用随机梯度类方法训练时，影响不大。

最终，得到的过程需要 $O(D)$ 的开销，再加上运行伴随灵敏度法（一种特定的数值求解方法）的成本。总的来说还不错，但远算不上出色。这里甚至不给出代码，因为扩大 CNF 的规模是个已知难题。那有没有其他选择？能做得更好吗？当然，好奇的读者，当然可以！

\subsection{顺流而行：流匹配}

\subsubsection{基本思想}

再看一遍前面介绍的 ODE。
\begin{equation}
  \frac{\mathrm d\mathbf x_t}{\mathrm dt}=v(\mathbf x_t,t).
  \tag{9.59}\label{eq:09-59}
\end{equation}
除了这个 ODE，再假设有一个\emph{已知}分布 $q_0(\mathbf x)$（例如标准高斯分布）和一个数据分布 $q_1(\mathbf x)$。

\begin{booknote}{旁注}
为与流匹配文献的惯例一致，现在从噪声（$t=0$）走向数据（$t=1$）。这一方向定义了正向动力学，也定义了生成过程。扩散模型和基于得分的模型采用相反的时间方向，从数据走向噪声。
\end{booknote}

从上面对 CNF 的讨论可知，任意时刻 $t$ 的分布由连续性方程刻画。利用瞬时变量变换，还可以求出一个解，即一个概率分布。不过，这些步骤听上去相当复杂，而且如前所述，训练的计算开销很大。那该怎么办？也可以换个问法。假如已知向量场 $v(\mathbf x_t,t)$ 和分布 $p_t(\mathbf x)$，能做些什么？我们该如何训练模型？模型又是什么？还记得得分匹配方法吗？再看看去噪得分匹配的损失。如果在这里采用类似方法，直接寻找向量场的模型 $v_\theta(\mathbf x_t,t)$ 呢？与得分匹配类似，我们求解如下形式的回归问题。
\begin{equation}
  \ell_{\mathit{FM}}(\theta)=
  \mathbb E_{t\sim U(0,1),\,\mathbf x_t\sim p_t(\mathbf x)}
  \bigl[\lVert v_\theta(\mathbf x_t,t)-v(\mathbf x_t,t)\rVert^2\bigr],
  \tag{9.60}\label{eq:09-60}
\end{equation}
这样就避开了 CNF 中寻找分布的过程。简单说，对任意均匀随机采样的时刻 $t$，从分布 $p_t(\mathbf x)$（假定已知！）采样 $\mathbf x_t$，使模型 $v_\theta(\mathbf x_t,t)$ 与\emph{真实}向量场 $v(\mathbf x_t,t)$（也假定已知！）之间的差异尽可能小。这个目标称为\textbf{流匹配}（FM）。

\dgmfig{fig_09_07}{.91\linewidth}{数据点（橙色的两弯月牙）周围，向量场模型（蓝色箭头）随时间变化的示例。}{fig:09-07}

为什么这样能奏效？理由很简单。如果 $\ell_{\mathit{FM}}(\theta)=0$，即模型完美模仿真实向量场，我们就能把任意噪声分布变换为数据分布！为什么？因为向量场会随着时间把各点\emph{推向}数据分布。看看\figref{fig:09-07}，蓝色箭头（向量场）表明，各点应如何随时间从噪声分布 $q_0(\mathbf x)$（例如标准高斯分布）\emph{演变}到数据分布 $q_1(\mathbf x)$（图中橙色月牙）。

第二个问题是，这个方法的好处在哪里？答案同样简单。这是一个回归问题，用的是均方误差损失！过程很直接，损失是性质良好的凸函数。可以用自动求导，并用任意深度学习库实现它。太好了！

不过，问题来了。我们并不知道 $p_t(\mathbf x)$ 和 $v(\mathbf x_t,t)$。好奇的读者，我知道，前面的讨论里你一直没忘这一点。你很客气，耐心等到了现在。换作我，早就说“胡扯”了。假设向量场和分布 $p_t(\mathbf x)$ 都已知，简直是胡扯。那该怎么办，你问？下面就来说明怎样处理！

\subsubsection{条件流匹配}

先考虑一个改写后的问题。引入从给定分布 $q(\mathbf z)$ 采样的附加变量 $\mathbf z$。\emph{条件} ODE 具有以下形式。
\begin{equation}
  \frac{\mathrm d\mathbf x_t}{\mathrm dt}=v(\mathbf x_t,t;\mathbf z).
  \tag{9.61}\label{eq:09-61}
\end{equation}
目前，可以把这个问题看作前面无条件 ODE 问题的替代。一般来说，只要条件信息相关，条件问题通常更容易处理。$\mathbf z$ 可以看成额外信息，例如数据 $\mathbf x_1$，也可以是类别标签、一段文本、一个音频信号或一幅附加图像。由于还需要从某个分布 $q(\mathbf z)$ 采样 $\mathbf z$，\textbf{条件流匹配}（conditional flow matching，CFM）的损失可以定义为
\begin{equation}
  \begin{aligned}
    \ell_{\mathit{CFM}}(\theta)
      ={}&\mathbb E_{t\sim U(0,1),\,\mathbf x_t\sim p_t(\mathbf x\mid\mathbf z),\,\mathbf z\sim q(\mathbf z)}
      \\[-1pt]
      &\bigl[\lVert v_\theta(\mathbf x_t,t)-v(\mathbf x_t,t;\mathbf z)\rVert^2\bigr],
  \end{aligned}
  \tag{9.62}\label{eq:09-62}
\end{equation}
这里仍用无条件模型来拟合条件向量场。

在 CFM 损失中，要定义每个 $t$ 的条件分布 $p_t(\mathbf x\mid\mathbf z)$，真实向量场也以 $\mathbf z$ 为条件，即 $v(\mathbf x_t,t;\mathbf z)$。乍看之下，CFM 问题似乎对 FM 问题没有什么用。为什么加上条件会有助于学习一个应当在无条件情形下工作的模型？文献\citerefs{32,33}证明，两种损失只相差一个与 $\theta$ 无关的常数，所以梯度相等！

\begin{booknote}{定理 \citeref{32}}
\textbf{定理 1}\quad 若对所有 $\mathbf x\in\mathbb R^D$ 和 $t\in[0,1]$ 都有 $p_t(\mathbf x)>0$，则 $\ell_{\mathit{FM}}$ 和 $\ell_{\mathit{CFM}}$ 只相差一个与 $\theta$ 无关的常数。因此，
\[
  \nabla_\theta\ell_{\mathit{FM}}(\theta)=\nabla_\theta\ell_{\mathit{CFM}}(\theta).
\]
\hfill\(\square\)
\end{booknote}

这个结果真妙！它意味着，用条件版本的损失训练的模型 $v_\theta(\mathbf x,t)$ 本身是无条件的，并且与无条件流匹配问题的解一致。太好了！一个问题解决了，我们可以用 CFM 代替 FM。现在又有了另一个问题，条件 $\mathbf z$ 该取什么，它的分布 $q(\mathbf z)$ 又是什么？好在有多种选择（例如参见文献\citeref{33}），这里集中介绍其中两种。
\begin{enumerate}
  \item 在文献\citeref{32}的 CNF 中，$\mathbf z$ 是一个数据点 $\mathbf x_1$，所以 $q(\mathbf z)=q_1(\mathbf z)$。换句话说，$q(\mathbf z)$ 是数据分布。
  \item 在文献\citeref{33}的 CNF 中，也就是独立条件流匹配（independent CFM，iCFM），$\mathbf z$ 是一对噪声与数据 $\mathbf z=(\mathbf x_0,\mathbf x_1)$，二者独立采样，即 $q_0(\mathbf z)=q(\mathbf x_0)q_1(\mathbf x_1)$。\footnote{译注：此处保留原书的下标写法。下文第~9.4.2.3~节对 iCFM 的定义为 $q(\mathbf z)=q_0(\mathbf x_0)q_1(\mathbf x_1)$。}
\end{enumerate}
iCFM 的一种扩展是通过求解最优传输（optimal transport）问题来采样 $\mathbf z=(\mathbf x_0,\mathbf x_1)$。感兴趣的读者（没错，就是你！）可以参阅文献\citeref{33}了解详情。

\subsubsection{条件概率路径}

现在知道，条件流匹配问题可以作为无条件流匹配问题的理想替代。不过，最后一块拼图是怎样获得条件分布 $p_t(\mathbf x\mid\mathbf z)$，即所谓的（条件）概率路径（conditional probability path）。从前面对 CNF 的讨论可知，可以通过连续性方程计算概率路径，但为此需要知道向量场。这是一个典型的循环论证。怎样避开它？

好奇的读者，你已经了解了不少人工智能知识，可能已经猜到这里能用什么最简单的方法了。我没听错吧，你在小声说“高斯分布”？没错！我们先考虑 $p_t(\mathbf x\mid\mathbf z)$ 的形式，再由此推导向量场 $v(\mathbf x,t;\mathbf z)$。这或许会限制整个方法的适用范围，可我们总得算出点什么，还有比高斯分布更好的选择吗？来看看下面形式的条件概率路径。
\begin{equation}
  p_t(\mathbf x\mid\mathbf z)=
  \mathcal N(\mathbf x\mid\mu(\mathbf z,t),\sigma^2(\mathbf z,t)\mathbf I),
  \tag{9.63}\label{eq:09-63}
\end{equation}
这是一个高斯分布，均值函数为 $\mu(\mathbf z,t)$，协方差矩阵为对角矩阵，标准差函数为 $\sigma(\mathbf z,t)$。一般来说，生成这些分布的 ODE 并不唯一。然而，下面的定理表明，可以由唯一的向量场得到这些分布！\footnote{译注：原书称“唯一的向量场”。式 \eqref{eq:09-64} 给出了一个仿射向量场。若在维数至少为 2 时不限制向量场的形式，可以加入绕均值旋转的速度场而保持相同的各向同性高斯路径。}

\begin{booknote}{定理 \citeref{32}}
\textbf{定理 2}\quad 在初始条件 $p_0(\mathbf x)=\mathcal N(\mu_0,\sigma_0^2\mathbf I)$ 下，产生
$p_t(\mathbf x\mid\mathbf z)=\mathcal N(\mathbf x\mid\mu(\mathbf z,t),\sigma^2(\mathbf z,t)\mathbf I)$
的唯一向量场具有以下形式。
\begin{equation}
  v(\mathbf x,t;\mathbf z)
  =\frac{\sigma'(\mathbf z,t)}{\sigma(\mathbf z,t)}
  \bigl(\mathbf x-\mu(\mathbf z,t)\bigr)+\mu'(\mathbf z,t),
  \tag{9.64}\label{eq:09-64}
\end{equation}
其中 $\sigma'(\mathbf z,t)$ 和 $\mu'(\mathbf z,t)$ 分别表示 $\sigma(\mathbf z,t)$ 和 $\mu(\mathbf z,t)$ 对时间的导数。
\end{booknote}

可以停下来想一想，或者再读一遍。明白了吗？没错，这个定理给了我们一个不可思议的结果！如果考虑高斯形式的一类条件概率路径，只要均值和标准差可微，就能\textbf{解析地}计算条件向量场。换句话说，我们有了流匹配的一套通用做法。第一，把均值和标准差定义为 $\mathbf z$ 与时间 $t$ 的函数。第二，对给定的均值和标准差函数形式，计算条件向量场。第三，针对给定的向量场模型，优化条件流匹配损失。就这样，非常简单。为了更好地理解它怎样工作，我们来看两种具体的条件流匹配形式。
\begin{itemize}
  \item \textbf{Lipman 等人的 CFM。} 后续代码中记为 \texttt{fm}。令 $\mathbf z\equiv\mathbf x_1$ 为从数据分布采样的数据，$\mathbf x_1\sim q_1(\mathbf x)$。然后将均值和标准差函数定义如下。
  \begin{align}
    \mu(\mathbf z,t)&=t\mathbf x_1,
      \tag{9.65}\label{eq:09-65}\\
    \sigma(\mathbf z,t)&=t\sigma_{\mathit{const}}-t+1,
      \tag{9.66}\label{eq:09-66}
  \end{align}
  其中 $\sigma_{\mathit{const}}>0$ 为平滑常数。因此，得到以下条件概率路径和条件向量场。
  \begin{align}
    p_t(\mathbf x\mid\mathbf z)
      &=\mathcal N\bigl(\mathbf x\mid t\mathbf x_1,(t\sigma_{\mathit{const}}-t+1)^2\mathbf I\bigr),
      \tag{9.67}\label{eq:09-67}\\
    v(\mathbf x,t;\mathbf z)
      &=\frac{\mathbf x_1-(1-\sigma_{\mathit{const}})\mathbf x}
        {1-(1-\sigma_{\mathit{const}})t}.
      \tag{9.68}\label{eq:09-68}
  \end{align}
  要求 $v(\mathbf x,t;\mathbf z)$ 的解析形式，需要应用上面的定理。这样得到一条概率路径，从标准高斯分布 $p_0(\mathbf x)=\mathcal N(\mathbf x\mid\mathbf 0,\mathbf I)$，到以某个数据点为中心、标准差为 $\sigma_{\mathit{const}}$ 的高斯分布 $p_1(\mathbf x)=\mathcal N(\mathbf x\mid\mathbf x_1,\sigma_{\mathit{const}}^2\mathbf I)$\citeref{32}。
  \item \textbf{Tong 等人的 iCFM。} 令 $\mathbf z\equiv(\mathbf x_0,\mathbf x_1)$，且 $q(\mathbf z)=q_0(\mathbf x_0)q_1(\mathbf x_1)$。接下来选择以下均值和标准差。
  \begin{align}
    \mu(\mathbf z,t)&=t\mathbf x_1+(1-t)\mathbf x_0,
      \tag{9.69}\label{eq:09-69}\\
    \sigma(\mathbf z,t)&=\sigma_{\mathit{const}},
      \tag{9.70}\label{eq:09-70}
  \end{align}
  其中 $\sigma_{\mathit{const}}>0$ 为平滑常数。因为 $t\in[0,1]$，均值函数是噪声与数据之间的\emph{插值}。得到的条件概率路径和条件向量场如下。
  \begin{align}
    p_t(\mathbf x\mid\mathbf z)
      &=\mathcal N\bigl(\mathbf x\mid t\mathbf x_1+(1-t)\mathbf x_0,
        \sigma_{\mathit{const}}^2\mathbf I\bigr),
      \tag{9.71}\label{eq:09-71}\\
    v(\mathbf x,t;\mathbf z)&=\mathbf x_1-\mathbf x_0.
      \tag{9.72}\label{eq:09-72}
  \end{align}
  有趣的是，得到的向量场是数据点与所采样噪声的差。这是因为假设概率路径上的标准差固定，所以应用定理 2 后只剩均值函数的导数。Tong 等人\citeref{33}表明（见该文的命题 3.3），边界分布为 $p_0(\mathbf x)=q_0(\mathbf x)*\mathcal N(\mathbf x\mid\mathbf 0,\sigma_{\mathit{const}}^2\mathbf I)$ 和 $p_1(\mathbf x)=q_1(\mathbf x)*\mathcal N(\mathbf x\mid\mathbf 0,\sigma_{\mathit{const}}^2\mathbf I)$，其中 $*$ 表示卷积算子。例如，若取 $q_0(\mathbf x)=\mathcal N(\mathbf x\mid\mathbf 0,\mathbf I)$，则 $p_0(\mathbf x)=\mathcal N(\mathbf x\mid\mathbf 0,(\sigma_{\mathit{const}}^2+1)\mathbf I)$。不过，$q_0(\mathbf x)$ 可以是任意分布，高斯分布只是一个例子。若 $q_1(\mathbf x)$ 是数据分布，另一个边界分布为 $p_1(\mathbf x)=\mathcal N(\mathbf x\mid\mathbf x_1,\sigma_{\mathit{const}}^2\mathbf I)$，与 Lipman 等人的 CFM 相同。\footnote{译注：给定数据点 $\mathbf x_1$ 时，终点的条件分布是这里的单个高斯分布。对 $\mathbf x_1$ 边缘化后，终点分布是前面给出的 $q_1$ 与高斯分布的卷积。}
\end{itemize}

\dgmfig{fig_09_08}{.95\linewidth}{流匹配方法的工作原理示意图。（a）Lipman 等人的 CNF 示例。（b）Tong 等人的 iCNF 示例。点线表示噪声与数据之间的插值。}{fig:09-08}[图中 time、noise distribution 和 data distribution 分别表示时间、噪声分布和数据分布。]

Lipman 等人的 CFM 与 Tong 等人的 iCFM 之间差异相当微妙。不过，这些细微差异会导致概率路径表现不同。\figref{fig:09-08} 展示了这两种 CFM 的示例。Lipman 等人的 CFM 从标准高斯分布出发，随时间演变为数据空间中较窄的高斯分布，即标准差减小，见\figref{fig:09-08}（a）。Tong 等人的 CFM 则定义一个较窄的高斯分布，并使它随时间\emph{移动}到数据空间，见\figref{fig:09-08}（b）。

\subsubsection{训练}

定义向量场模型的训练所需的组件已经齐备。通用算法如下。
\begin{booknote}[colframe=blue!65!black,colback=blue!4]{流匹配的训练}
  \begin{enumerate}
    \item 采样 $t\sim\operatorname{Uniform}(0,1)$。
    \item 采样 $\mathbf z\sim q(\mathbf z)$。
    \item 计算 $\mu_t(\mathbf z)$ 和 $\sigma_t(\mathbf z)$。
    \item 采样 $\mathbf x_t\sim\mathcal N(\mathbf x\mid\mu(\mathbf z,t),\sigma(\mathbf z,t))$。
    \item 计算向量场 $v(\mathbf x_t,t;\mathbf z)$。
    \item 计算损失 $\ell_{\mathit{CFM}}(\theta)=\lVert v_\theta(\mathbf x_t,t)-v(\mathbf x_t,t;\mathbf z)\rVert^2$。
    \item 更新参数 $\theta\gets\operatorname{Update}(\theta,\nabla_\theta\ell_{\mathit{CFM}}(\theta))$。
  \end{enumerate}
\end{booknote}
选择不同的均值函数和标准差函数，会得到特定形式的向量场。最终，第 3 步和第 5 步得到相应的具体公式，其余步骤相同。

\subsubsection{从流匹配模型采样}

好奇的读者，再提醒一句，我们学到的是向量场的模型。有了这个模型，采样就很直接。
\begin{booknote}[colframe=blue!65!black,colback=blue!4]{采样}
  \begin{enumerate}
    \item 采样 $\mathbf x\sim q_0(\mathbf x)$。
    \item 以 $\Delta$ 为步长运行正向欧拉法，直到 $t=1$。
    \[
      \mathbf x_{t+\Delta}=\mathbf x_t+v_\theta(\mathbf x_t,t)\Delta.
    \]
  \end{enumerate}
\end{booknote}
请注意，流匹配假定时间从 $t=0$（即噪声）走向 $t=1$（即数据），与基于得分的生成模型所取的方向相反。因此，向量场模型训练完成后，采用正向欧拉法。基于得分的生成模型则采用反向欧拉法。

\subsubsection{计算对数似然函数}

讨论 CNF 时，我们说明了如何计算（对数）似然函数。流匹配训练时不需要对数似然函数，所以至今没有考虑它。然而，能够计算对数似然函数值，对评估尤其有用，在许多实际应用中也有好处。流匹配对数似然函数的完整推导见文献\citeref{32}附录 C。主要思想与 CNF 类似。为了使计算切实可行，采用 Hutchinson 迹估计量，得到
\begin{equation}
  \ln p_1(\hat{\mathbf x}_1)\approx\ln p_0(\hat{\mathbf x}_0)-f_0,
  \tag{9.73}\label{eq:09-73}
\end{equation}
其中 $\hat{\mathbf x}_1$ 是一个数据点，$\hat{\mathbf x}_0$ 是与 $\hat{\mathbf x}_1$ 对应的噪声，$f_0$ 是向量场雅可比矩阵的迹的近似。注意，给定的只有 $\hat{\mathbf x}_1$，但可以通过以下过程得到 $\hat{\mathbf x}_0$ 和 $f_0$\citeref{32}。
\begin{booknote}[colframe=blue!65!black,colback=blue!4]{计算流匹配模型的对数似然}
  \begin{enumerate}
    \item 对给定数据点 $\hat{\mathbf x}_1$，设置初始条件。
    \[
      \begin{bmatrix}\phi_1\\f_1\end{bmatrix}
      =\begin{bmatrix}\hat{\mathbf x}_1\\0\end{bmatrix}.
    \]
    \item 定义如下 ODE。
    \[
      \frac{\mathrm d}{\mathrm ds}
      \begin{bmatrix}\phi_{1-s}\\f_{1-s}\end{bmatrix}
      =\begin{bmatrix}
        -v_\theta(\phi_{1-s},1-s)\\
        \epsilon^\top\nabla_\phi v_\theta(\phi_{1-s},1-s)\epsilon
      \end{bmatrix}
    \]
    其中 $\epsilon\sim\mathcal N(\mathbf 0,\mathbf I)$。

    \emph{注意。} 用自动求导计算 $\nabla_\phi v_\theta(\phi_{1-s},1-s)$。
    \item 用反向欧拉法求解第 2 步的 ODE。

    \emph{注意。} 这里需要从 $t=1$ 回到 $t=0$，所以采用反向欧拉法。
    \item 输出结果。
    \[
      \begin{bmatrix}\phi_0\\f_0\end{bmatrix}
      =\begin{bmatrix}\hat{\mathbf x}_0\\\hat c\end{bmatrix}.
    \]
  \end{enumerate}
\end{booknote}
将这个过程的输出代入式 \eqref{eq:09-73}，得到
\begin{equation}
  \ln p_1(\hat{\mathbf x}_1)\approx\ln p_0(\hat{\mathbf x}_0)-\hat c.
  \tag{9.74}\label{eq:09-74}
\end{equation}
好奇的读者，请记住，由于使用了迹估计量，这是一个近似。不过，这个估计量是无偏的。对给定数据点运行这个过程 $M$ 次，再取平均，可以减小方差。

\subsection{计算对数似然函数}

我们已经概述了如何用流匹配实现训练和采样（生成）。代码相当直接，见下面的代码片段。这个实现选择标准高斯噪声分布，但请记住，其他分布也可以。

\Needspace{6\baselineskip}
\bookcode{code/ch09_03_flow_matching.py}{条件流匹配类。}{lst:09-03}

在展示生成图像之前，先看看流匹配在双月牙数据集上的表现，见\figref{fig:09-09}。从按标准高斯分布生成的数据出发，向量场模型将各点推向两弯橙色月牙。\footnote{译注：代码清单~\ref{lst:09-03} 保留原书写法。其 \texttt{log\_prob} 函数逐步覆盖 \texttt{f\_t}，未累加各步的迹估计。第 135 行的注释说切回训练模式，但该行仍调用 \texttt{eval()}。}

\dgmfig{fig_09_09}{.91\linewidth}{用流匹配训练出的向量场模型采样的示例。采样点（黑色叉号）向数据（橙色月牙）演变。}{fig:09-09}

\dgmfig{fig_09_10}{\linewidth}{部分结果。（a）真实图像的样本。（b）生成图像的样本。（c）验证集上计算的得分匹配损失示例。}{fig:09-10}[（c）横轴 epochs 为训练轮次，纵轴 nll 为负对数似然。原书图注称“得分匹配损失”，正文说明报告的是估计的负对数似然。]

向量场模型采用 MLP，超参数设为 $\sigma_{\mathit{const}}=0.1$ 和 $T=100$。运行代码后，可以得到类似\figref{fig:09-10} 的结果。注意，这里报告的是按前述过程估计的负对数似然。

\subsection{流匹配的未来是什么？}

流匹配这个框架最早见于 2022 年秋发表的文献\citeref{32}。此后已经有人提出多种扩展。这里我会指出其中几种，不过，好奇的读者，我强烈建议你自己查一查新的进展，因为我觉得这个框架格外有意思！请留意以下论文。
\begin{itemize}
  \item 与潜在扩散类似，文献\citeref{34}在自编码器设定中将流匹配用作“先验”。先训练自编码器，再在潜在空间中训练向量场模型。之后，从 FM 模型中采样，将样本解码回数据空间。
  \item CFM 中的插值思想进一步扩展为文献\citerefs{35,36}提出的\emph{随机插值量}（stochastic interpolants）。我强烈建议查阅这些论文，它们给出了许多有趣的扩展，既有理论方面的，也有实践方面的。
  \item 如前所述，可以用最优传输（OT）提出更好的分布 $q(\mathbf z)$。由此得到 OT-CFM\citeref{33}和薛定谔桥（Schrödinger bridge）CFM\citerefs{29,37}。
  \item 文献\citeref{38}提出的作用量匹配（action matching）与 CFM 密切相关。它可以学习各点随时间移动的底层机制，省去对每一步分布的建模。
  \item 这里考虑的是欧氏空间中的插值。Chen 和 Lipman\citeref{39}提出了将 CFM 扩展到一般黎曼流形的方法。
  \item 为利用数据的对称性，文献\citeref{40}修改了 OT-CFM 中的代价函数，将这些对称性考虑进去。研究者还使用等变图神经网络，构造了等变流匹配（equivariant flow matching）。
  \item 这里讨论了连续随机变量的情形。不过，在实践中经常要处理离散数据，例如分子、蛋白质和像素值。Campbell 等人\citeref{41}提出了离散流模型（discrete flow models），可将其看作处理离散数据的 CFM 版本。
  \item 还有研究把 CFM 用于基于仿真的推断（simulation-based inference）\citeref{42}。这类问题允许访问模拟器，但似然函数未知或难以计算。
  \item 基于得分的生成模型与流匹配联系非常紧密。要了解更多细节，我强烈推荐文献\citerefs{33,5}，后一篇的附录实在精彩！
\end{itemize}

还有许多优秀的论文没能提到。这并非有意遗漏，要跟上所有研究实在太难了。这也表明，流匹配正在成为生成式人工智能的主流研究方向。依我看，这是一个迷人的框架，可能带来多项突破！


\FloatBarrier
% 原书印刷页 254–257，附件 PDF 第 269–272 页。
\begin{chapterreferences}
\refentry{1} Yang Song, Jascha Sohl-Dickstein, Diederik P Kingma, Abhishek Kumar, Stefano Ermon, and Ben Poole. Score-based generative modeling through stochastic differential equations. In \emph{International Conference on Learning Representations}, 2020.
\refentry{2} Jascha Sohl-Dickstein, Eric Weiss, Niru Maheswaranathan, and Surya Ganguli. Deep unsupervised learning using nonequilibrium thermodynamics. In \emph{International Conference on Machine Learning}, pages 2256–2265. PMLR, 2015.
\refentry{3} Jonathan Ho, Ajay Jain, and Pieter Abbeel. Denoising diffusion probabilistic models. \emph{Advances in Neural Information Processing Systems}, 33:6840–6851, 2020.
\refentry{4} Diederik P Kingma, Tim Salimans, Ben Poole, and Jonathan Ho. Variational diffusion models. \emph{arXiv preprint arXiv:2107.00630}, 2021.
\refentry{5} Diederik Kingma and Ruiqi Gao. Understanding diffusion objectives as the elbo with simple data augmentation. \emph{Advances in Neural Information Processing Systems}, 36, 2024.
\refentry{6} Max Welling and Yee W Teh. Bayesian learning via stochastic gradient langevin dynamics. In \emph{Proceedings of the 28th international conference on machine learning (ICML-11)}, pages 681–688. Citeseer, 2011.
\refentry{7} Aapo Hyvärinen and Peter Dayan. Estimation of non-normalized statistical models by score matching. \emph{Journal of Machine Learning Research}, 6(4), 2005.
\refentry{8} Aapo Hyvärinen. Some extensions of score matching. \emph{Computational statistics \& data analysis}, 51(5):2499–2512, 2007.
\refentry{9} Pascal Vincent. A connection between score matching and denoising autoencoders. \emph{Neural computation}, 23(7):1661–1674, 2011.
\refentry{10} Yang Song and Stefano Ermon. Generative modeling by estimating gradients of the data distribution. \emph{arXiv preprint arXiv:1907.05600}, 2019.
\refentry{11} Yang Song and Stefano Ermon. Improved techniques for training score-based generative models. \emph{Advances in neural information processing systems}, 33:12438–12448, 2020.
\refentry{12} Yang Song, Sahaj Garg, Jiaxin Shi, and Stefano Ermon. Sliced score matching: A scalable approach to density and score estimation. In \emph{Uncertainty in Artificial Intelligence}, pages 574–584. PMLR, 2020.
\refentry{13} Julien Rabin, Gabriel Peyré, Julie Delon, and Marc Bernot. Wasserstein barycenter and its application to texture mixing. In \emph{Scale Space and Variational Methods in Computer Vision: Third International Conference, SSVM 2011, Ein-Gedi, Israel, May 29–June 2, 2011, Revised Selected Papers 3}, pages 435–446. Springer, 2012.
\refentry{14} Benjamin Marlin, Kevin Swersky, Bo Chen, and Nando Freitas. Inductive principles for restricted boltzmann machine learning. In \emph{Proceedings of the thirteenth International Conference on Artificial Intelligence and Statistics}, pages 509–516, 2010.
\refentry{15} Kevin Swersky, Marc’Aurelio Ranzato, David Buchman, Nando D Freitas, and Benjamin M Marlin. On autoencoders and score matching for energy based models. In \emph{Proceedings of the 28th international conference on machine learning (ICML-11)}, pages 1201–1208, 2011.
\refentry{16} Zengyi Li, Yubei Chen, and Friedrich T Sommer. Learning energy-based models in high-dimensional spaces with multiscale denoising-score matching. \emph{Entropy}, 25(10):1367, 2023.
\refentry{17} Chenlin Meng, Kristy Choi, Jiaming Song, and Stefano Ermon. Concrete score matching: Generalized score matching for discrete data. \emph{Advances in Neural Information Processing Systems}, 35:34532–34545, 2022.
\refentry{18} Shitong Luo and Wei Hu. Score-based point cloud denoising. In \emph{Proceedings of the IEEE/CVF International Conference on Computer Vision}, pages 4583–4592, 2021.
\refentry{19} Simo Särkkä and Arno Solin. \emph{Applied stochastic differential equations}, volume 10. Cambridge University Press, 2019.
\refentry{20} Yang Song, Conor Durkan, Iain Murray, and Stefano Ermon. Maximum likelihood training of score-based diffusion models. \emph{Advances in neural information processing systems}, 34:1415–1428, 2021.
\refentry{21} Yang Song and Diederik P Kingma. How to train your energy-based models. \emph{arXiv preprint arXiv:2101.03288}, 2021.
\refentry{22} Cheng Lu, Yuhao Zhou, Fan Bao, Jianfei Chen, Chongxuan Li, and Jun Zhu. Dpm-solver: A fast ode solver for diffusion probabilistic model sampling in around 10 steps. \emph{Advances in Neural Information Processing Systems}, 35:5775–5787, 2022.
\refentry{23} Zhenyu Zhou, Defang Chen, Can Wang, and Chun Chen. Fast ode-based sampling for diffusion models in around 5 steps. \emph{arXiv preprint arXiv:2312.00094}, 2023.
\refentry{24} Tero Karras, Miika Aittala, Timo Aila, and Samuli Laine. Elucidating the design space of diffusion-based generative models. \emph{Advances in Neural Information Processing Systems}, 35:26565–26577, 2022.
\refentry{25} Arash Vahdat, Karsten Kreis, and Jan Kautz. Score-based generative modeling in latent space. \emph{arXiv preprint arXiv:2106.05931}, 2021.
\refentry{26} Ricky TQ Chen, Yulia Rubanova, Jesse Bettencourt, and David K Duvenaud. Neural ordinary differential equations. \emph{Advances in neural information processing systems}, 31, 2018.
\refentry{27} Kaiwen Zheng, Cheng Lu, Jianfei Chen, and Jun Zhu. Improved techniques for maximum likelihood estimation for diffusion odes. In \emph{International Conference on Machine Learning}, pages 42363–42389. PMLR, 2023.
\refentry{28} Yang Song, Prafulla Dhariwal, Mark Chen, and Ilya Sutskever. Consistency models. \emph{arXiv preprint arXiv:2303.01469}, 2023.
\refentry{29} Valentin De Bortoli, James Thornton, Jeremy Heng, and Arnaud Doucet. Diffusion schrödinger bridge with applications to score-based generative modeling. \emph{Advances in Neural Information Processing Systems}, 34:17695–17709, 2021.
\refentry{30} Heli Ben-Hamu, Samuel Cohen, Joey Bose, Brandon Amos, Aditya Grover, Maximilian Nickel, Ricky TQ Chen, and Yaron Lipman. Matching normalizing flows and probability paths on manifolds. \emph{arXiv preprint arXiv:2207.04711}, 2022.
\refentry{31} Will Grathwohl, Ricky TQ Chen, Jesse Bettencourt, Ilya Sutskever, and David Duvenaud. Ffjord: Free-form continuous dynamics for scalable reversible generative models. \emph{arXiv preprint arXiv:1810.01367}, 2018.
\refentry{32} Yaron Lipman, Ricky TQ Chen, Heli Ben-Hamu, Maximilian Nickel, and Matt Le. Flow matching for generative modeling. \emph{arXiv preprint arXiv:2210.02747}, 2022.
\refentry{33} Alexander Tong, Nikolay Malkin, Guillaume Huguet, Yanlei Zhang, Jarrid Rector-Brooks, Kilian Fatras, Guy Wolf, and Yoshua Bengio. Improving and generalizing flow-based generative models with minibatch optimal transport. \emph{arXiv preprint arXiv:2302.00482}, 2023.
\refentry{34} Quan Dao, Hao Phung, Binh Nguyen, and Anh Tran. Flow matching in latent space. \emph{arXiv preprint arXiv:2307.08698}, 2023.
\refentry{35} Michael S Albergo, Nicholas M Boffi, and Eric Vanden-Eijnden. Stochastic interpolants: A unifying framework for flows and diffusions. \emph{arXiv preprint arXiv:2303.08797}, 2023.
\refentry{36} Michael S Albergo, Mark Goldstein, Nicholas M Boffi, Rajesh Ranganath, and Eric Vanden-Eijnden. Stochastic interpolants with data-dependent couplings. \emph{arXiv preprint arXiv:2310.03725}, 2023.
\refentry{37} Francisco Vargas, Pierre Thodoroff, Austen Lamacraft, and Neil Lawrence. Solving schrödinger bridges via maximum likelihood. \emph{Entropy}, 23(9):1134, 2021.
\refentry{38} Kirill Neklyudov, Daniel Severo, and Alireza Makhzani. Action matching: A variational method for learning stochastic dynamics from samples, 2022. ICML, 6662.
\refentry{39} Ricky TQ Chen and Yaron Lipman. Riemannian flow matching on general geometries. \emph{arXiv preprint arXiv:2302.03660}, 2023.
\refentry{40} Leon Klein, Andreas Krämer, and Frank Noé. Equivariant flow matching. \emph{Advances in Neural Information Processing Systems}, 36, 2024.
\refentry{41} Andrew Campbell, Jason Yim, Regina Barzilay, Tom Rainforth, and Tommi Jaakkola. Generative flows on discrete state-spaces: Enabling multimodal flows with applications to protein co-design. \emph{arXiv preprint arXiv:2402.04997}, 2024.
\refentry{42} Jonas Wildberger, Maximilian Dax, Simon Buchholz, Stephen Green, Jakob H Macke, and Bernhard Schölkopf. Flow matching for scalable simulation-based inference. \emph{Advances in Neural Information Processing Systems}, 36, 2024.
\end{chapterreferences}

